\section{背景与相关工作}
\label{sec:related}

Next连接基于证书的连续续花、备付支撑的公共执行及授权历史修复。
最接近的工作已经在认证、抵押和受控编辑方面建立先例。
我们依据缺失资金来源的责任归属，以及责任落实过程中保持的付款与区块身份，界定Next的研究定位。

\subsection{基于证书的支付与执行}

FastPay通过拜占庭一致广播实现低延迟支付，
并支持由预充值资金支撑的主账本部署~\cite{fastpay}。
其可兑付性论证允许在相关入账转账尚未全部在本地确认时确认有效证书，
前提是补齐发送方先前转出付款的缺失证书。
兑付提供了一条基于证书从主账本资金池获取资金的路径。
Zef将可弃置账户与具有隐私性的可转让币相结合。
正常消费使用当前币证书及消费证据，而非完整的祖先交易历史~\cite{zef}。

Sui Lutris对归属特定所有者的对象采用广播执行，对共享对象的交互采用共识~\cite{lutris}。
Stingray通过拜占庭有界计数器和FastUnlock支持并发交易~\cite{stingray}；
其重叠额度与回收机制是与Next资源核算密切相关的先例。
其可交换事务工作负载所处理的争用问题，与每笔交易消费前序输出的支付链有所不同。
Next关注在资金来源交易执行之前，公共执行对经认证UTXO的消费，
同时登记由直接发行组织承担、具有资金保障的直接义务。
全部输出金额的预留与累计损失的保留，既计入已经公共登记的责任，
也计入隐藏证书所代表的承诺。

\subsection{通道与抵押担保}

Blitz使用pay-or-revoke替代通道支付中的两阶段提交，
Donner则构建应对Domino攻击的虚拟通道~\cite{blitz,donner}。
CryptoMaze提供部分付款不可链接的原子多路径通道支付，
同时避免在共享边上重复建立链下合约~\cite{cryptomaze}。
这些机制基于通道安排运行。
Next使用具有已注册消费路由和组织备付资金的原生账本，
因此交易双方既无需双边通道，也无需在彼此之间预先建立注入流动性的通道路径。

Snappy通过抵押支持的批准机制保护接受快速链上付款的商户。
可复用支付担保将隐私保护与可在规定时间窗口内重复使用的客户抵押相结合~\cite{snappy,reusable}。
Next则将用户付款本金与组织提供的额外担保资金分离。
其赔付关闭已执行子付款的缺失资金来源，而不再次向该子付款的收款人增加本金。

HyperService通过公共Network Status Blockchain与Insurance Smart Contract，
对跨链执行违约进行仲裁，并在资金层面撤销已提交交易的效果~\cite{web3}。
Next处理的是单个原生UTXO账本内不同的责任边界：
赔付结清发行组织的直接义务，同时保持已执行后继付款的本金效果。

\subsection{受控编辑与历史一致性}

变色龙哈希允许在同一摘要下存在经授权的替代消息，并为可编辑区块链设计提供基础~\cite{redactable}。
SDR研究去中心化的基于属性的编辑权限与动态委托；
Escaping将委员会选择和投票移至共识层之外，以减少修订等待~\cite{sdr,escaping}。
Reconstructing Chameleon Hash提出更强的不可区分性与抗碰撞性定义~\cite{reconstructCH}。
Next采用已有适配原语，由公共支付状态检查提供实施兼容替换的授权。

Reparo提供公开可验证的修复层，并可在修复状态时重新执行受影响的交易~\cite{reparo}。
Next固定付款授权和输出：公共修复命令对赔付进行核算，
历史物化则替换获准修改的资金来源表示，而不重新执行后继付款。
原始命令和后续修复仍构成经济重放历史。

RedactChain协调同块并发编辑冲突与预处理~\cite{redactchain}。
CDEdit使用可复用权限令牌，减少多次编辑中重复的令牌和密钥生成及验证工作~\cite{cdedit}。
Next的批处理在共享最终区块表示的同时，保留每项直接义务的资格、备付扣款和终态检查。
其完整区块身份由稳定区块头哈希与原始分片集头共同构成。
经公共提交的逻辑修订授权确切的替换正文；物理安装可直接推进到后续已授权修订。
这一设计将仅生效一次的经济赔付与表示更新的物理安装分离，
同时保持后继引用，并支持重放同一公共命令历史。
